\documentclass[11pt]{article}
\usepackage[a4paper,margin=27mm]{geometry}
\usepackage{amsmath,amssymb,amsthm,booktabs,longtable,hyperref}
\usepackage{xurl}
\newtheorem{theorem}{Theorem}
\newtheorem{lemma}[theorem]{Lemma}
\newtheorem{proposition}[theorem]{Proposition}
\newtheorem{corollary}[theorem]{Corollary}
\newcommand{\B}{\mathrm B}
\newcommand{\Sph}{\mathrm S}
\newcommand{\note}[1]{\nolinkurl{#1}}
\title{A six-layer bound for quartic graphs of girth five\\and the minimum order at diameter twenty-two}
\author{Computational reconstruction of a girth--diameter bound}
\date{}
\begin{document}
\maketitle
\begin{abstract}
We prove that six consecutive nonsingleton distance layers of a simple
quartic graph of girth at least five contain at least twenty-two vertices.
The last local cases are excluded by twenty-seven explicit Boolean
refutations, whose formulas are reconstructed and whose proofs are checked
without a satisfiability solver. Combining this local bound with boundary
balls and a mixed-radius packing inequality shows that a quartic graph of
girth at least five with a vertex of eccentricity twenty-two has at least
ninety-five vertices. An included ninety-five-vertex graph has degree four,
girth five and diameter twenty-two. Thus $n(4;5,22)=95$. No claim about the
number of nonisomorphic minimum graphs is needed.
\end{abstract}

\section{Statement, notation and antecedents}

For fixed degree $k$, girth $g$ and diameter $d$, write $n(k;g,d)$ for the
minimum order of a connected simple $k$-regular graph with those parameters.
Table~5 of Cambie, Goedgebeur, Jooken and Van den Eede~\cite{source}
records $n(4;5,22)=95$ and at least one such graph. We give a lower bound
for a slightly larger domain and separately verify a graph attaining it.

\begin{theorem}\label{thm:main}
Let $G$ be a connected simple $4$-regular graph of girth at least five.
If some vertex has eccentricity $22$, then $|V(G)|\geq95$.
\end{theorem}

Fix a root $o$ and write $V_i=\{x:d_G(o,x)=i\}$. The symbols $V_i$ always
refer to these original distance layers. A ball taken in an induced
subgraph $H$ is written $\B_r^H(x)$, and its sphere is $\Sph_r^H(x)$.
Every edge changes the original index by at most one. Consequently a ball
of radius $r$ centred at original index $j$ is supported on indices
$j-r,\ldots,j+r$, even when the ball is taken in an induced subgraph.
We never identify reversed original layers with distance layers from a
new root.

The Moore-tree argument gives $|\B_2(x)|=1+4+12=17$ whenever all relevant
vertices have degree four: triangles and quadrangles are exactly the
obstructions to distinctness. Five consecutive nonempty original layers
also have mass at least $17$: take a geodesic through their middle layer
and its radius-two ball. The distance-layer framework and these Moore
interval bounds are reused from~\cite{source}. The elementary methods
carry no claim of priority. Detailed hand proofs used below are included
as companion notes, rather than replaced by a numerical search.

\section{The local six-layer theorem}

\begin{theorem}\label{thm:local}
Six consecutive original distance layers, each containing at least two
vertices, have total mass at least $22$.
\end{theorem}

The hand arguments in \note{premises/six_layer_mass_bound.rst} and
\note{premises/six_layer_twenty_equality_exclusion.rst} establish mass at
least $21$. We describe the remaining equality cases and the exact
interface to the finite certificates.

Suppose the six layers have total mass $21$, and number them locally
$0,\ldots,5$. They begin after the root layer, which is a singleton.
Choose a geodesic across them, with adjacent centres $u$ in layer $2$
and $v$ in layer $3$. Their other neighbours are $A_1,A_2,A_3$ and
$B_1,B_2,B_3$, labelled so that $A_1$ is in layer $1$ and $B_1$ in layer
$4$ on the geodesic. No $A$--$B$ edge is possible. A common second
neighbour $C_{ij}$ has parents $A_i,B_j$; at most one exists for each
cell $ij$, and it cannot have a second parent on either side. These
assertions follow respectively from the absence of triangles or
quadrangles. If $m$ grid cells are occupied, then
\[
 |\B_2(u)\cap\B_2(v)|=8+m,\qquad
 |\B_2(u)\cup\B_2(v)|=26-m.                         \tag{1}
\]
For each missing cell $ij$, name a private neighbour $U_{ij}$ of $A_i$
and a private neighbour $V_{ij}$ of $B_j$. All these private vertices
are distinct. Each of the eight centre or parent vertices is already
saturated by its four prescribed neighbours.

Cell $11$ is missing: a common neighbour would connect original indices
whose difference is three by a path of length two. Thus $m\leq8$.
Equation~(1) immediately settles $m\leq4$. At total mass $21$ the only
remaining possibilities are
\[
\begin{array}{c|rrrr}
m&8&7&6&5\\\hline
\text{vertices outside the displayed union}&3&2&1&0.
\end{array}
\]
The $m=8$ case is excluded in
\note{premises/six_layer_equality_eight_cells.rst}.
For $m=7$ the two missing cells either share a row or column, or lie in
distinct rows and columns. Their respective exclusions are in
\note{premises/six_layer_adjacent_holes_proof.rst} and
\note{premises/six_layer_equality_matching_holes.rst}.
The former is the corrected full proof, including its permitted
five-cycle subcase. These hand arguments do not assume that reversing
the interval preserves the original predecessor axiom.

\subsection{A necessary Boolean model for $m=5,6$}

For $m=5$ there are exactly twenty-one displayed vertices. For $m=6$ add
one vertex $X$. Private names can be chosen so that the six geodesic
vertices are
\[
 U_{11},A_1,u,v,B_1,V_{11},
\]
in local layers $0,1,2,3,4,5$, respectively. For each vertex introduce six
Boolean layer variables with exactly one true entry. The prescribed
edges are constants true; all other edges incident to the saturated
eight vertices are false. All remaining unordered pairs receive an edge
variable. The formula requires:
\begin{enumerate}
\item the six anchors above and at least two vertices in every layer;
\item endpoints of an edge to have indices differing by at most one;
\item degree at most four everywhere, and degree four at local indices
      $1,\ldots,4$;
\item degree at most three inside the interval at local index zero;
\item for each vertex at index $h>0$, a neighbour at index $h-1$;
\item at most one common neighbour for any pair, and none for an
      adjacent pair.
\end{enumerate}
There is no successor condition. At index zero there is no positive
internal-degree requirement: its compulsory predecessor is outside the
interval. Five-cycles remain allowed. The last condition excludes
quadrangles and triangles, not pentagons.

Every actual graph fragment satisfying our hypotheses gives a satisfying
assignment. Assign its original indices and its induced adjacency, then
evaluate the equivalence-defined Boolean gates. Interior degrees are
complete because an edge cannot leave the interval from an interior
index. At index zero the original nonroot layer supplies a predecessor
outside. All other requirements were just established. The deterministic
formula source and its Boolean gate library are included unchanged; only
the latter's elementary gates are used, not its unrelated graph builder.
A full encoding proof is \note{premises/proof.rst}.

\subsection{Complete coverage and checked refutations}

The available label symmetries are independent swaps of rows $2,3$ and
columns $2,3$. They fix cell $11$, all geodesic anchors and all layer
indices. They transport private names with their parents. Hence this
$C_2\times C_2$ action preserves every condition above. Neither transpose
nor reversal is used. There are
\[
 \binom82=28\quad(m=6),\qquad \binom83=56\quad(m=5)
\]
hole sets containing $11$, partitioned into $10$ and $17$ orbits.
\note{data/girth_diameter_bounds/hole_pattern_partition.json} lists every
member and every one of its four images; the companion coverage proof
justifies the action.

For each of the $27$ representatives, the package contains a deterministic
DIMACS formula, its primary-variable map and a complete proof stream.
The replay reconstructs the entire formula and map byte for byte, then
checks the saved proof with the included reverse-unit-propagation (RUP)
checker. Each added clause must be justified; an explicit empty-clause
derivation is required. The checker retains the fixed original-unit
consequences and a historical fallback, whose soundness arguments are
in \note{premises/root_units.tex} and
\note{premises/root_units_active.tex}. Its sources and imported source
are included; a local C++17 compiler rebuilds the executable. The
solver's original answer is not a premise, and no solver is required for
replay. The exact $27$ proofs contain $43\,029$ checked additions in total.

Thus each necessary formula is contradictory. The label action covers
every $m=5,6$ hole pattern, so all equality-$21$ cases are excluded.
Together with the hand lower bound $21$, this proves
Theorem~\ref{thm:local}.

\section{Boundary facts and singleton layers}

If an internal original layer is a singleton $\{w\}$, its two cuts are
positive even integers. Indeed every degree is even, both cuts are
nonempty, and their sum is four. They are therefore both two. The
induced prefix and suffix including $w$ are connected graphs in which
$w$ has degree two and every other vertex has degree four; call either
one a \emph{cap}. The induced section between consecutive singleton
layers has two degree-two boundary vertices and all other degrees four.
Connectivity follows from original root geodesics.

We use the following hand-proved facts from the included notes.
\begin{lemma}\label{lem:premises}
Under the preceding degree and girth hypotheses:
\begin{enumerate}
\item Every cap has at least $20$ vertices.
\item A boundary ball $\B_3(w)$ has at least $15$ vertices when no other
      exceptional boundary occurs nearby. A boundary ball $\B_4(w)$ has
      at least $19$ vertices if all its other vertices have degree four,
      the induced graph has even degrees, and a vertex at distance at
      least five exists.
\item If $|G|\leq94$ and $o$ has eccentricity $22$, the complete set of
      internal singleton indices is one of
      $\varnothing,\{4\},\{18\},\{4,18\}$.
\end{enumerate}
\end{lemma}
The cap proof is \note{premises/singleton_cap_twenty_bound.rst}; its
earlier degree-two cap argument is retained in
\note{premises/singleton_cap_bound.tex}.
The boundary proof is \note{premises/boundary_radius_four.rst}.
The position reduction uses
\note{premises/singleton_layer_section.tex},
\note{premises/singleton_pattern_bound.rst},
\note{premises/singleton_free_section_bounds.rst}, and the cap and
six-layer bounds above. The stages reduce the necessary nonempty
patterns to $66$, then $14$, then $12$, then $5$, and finally the three
displayed nonempty sets. The last step already follows from mass $21$,
in \note{premises/six_layer_twenty_equality_exclusion.rst}.
These are reductions of the complete set of singleton positions, not
selections of convenient individual singletons.

\section{A mixed-radius boundary inequality}

\begin{lemma}\label{lem:mixed}
Under the radius-four hypotheses of Lemma~\ref{lem:premises},
\[
 |\B_3(w)|+|\B_4(w)|\geq35.                           \tag{2}
\]
\end{lemma}
\begin{proof}
The first two spheres have sizes two and six. Split the six vertices
into triples $P,Q$ according to their first-sphere parents. Their
internal edges form a matching of size $m\leq3$. Every third-sphere
vertex has at most one parent in each triple, so $18-2m\geq12$ incidences
enter that sphere and $c=|\Sph_3(w)|\geq6$.

If $c\geq7$, then $|\B_3(w)|=9+c\geq16$, and the separate bound
$|\B_4(w)|\geq19$ proves (2). If $c=6$, the incidence count forces
$m=3$ and two parents for every third-sphere vertex. Label the three
matching edges $p_iq_i$, $i=0,1,2$. The six third-sphere labels are
exactly the off-diagonal pairs $ij$; they are independent by the
triangle and quadrangle restrictions with the parents. Thus twelve
incidences enter the fourth sphere.

A fourth-sphere vertex cannot meet two labels sharing a coordinate,
as that would give a quadrangle. It has at most three neighbours in
the third sphere. The only possible three-label sets are
$\{01,12,20\}$ and $\{10,02,21\}$. Neither set can repeat at two
fourth-sphere vertices, again by a quadrangle. Four vertices in the
fourth sphere could therefore receive at most $3+3+2+2=10$ incidences.
Its size is at least five. Hence $|\B_3(w)|=15$ and
$|\B_4(w)|\geq20$, proving (2).
\end{proof}
This is a joint inequality. It does not assert a universal boundary
radius-four lower bound of twenty.

\section{The global packing}

Assume for contradiction that $|G|\leq94$. Lemma~\ref{lem:premises}
gives four alternatives. A ball of radius two at $o$ has size $17$ and
support $0,1,2$. A radius-two ball at any vertex of $V_{22}$ also has
size $17$ and support $20,21,22$, even if $V_{22}$ is a singleton.

If there is no internal singleton, use those two end balls and the
disjoint intervals $3,\ldots,8$, $9,\ldots,14$, and $15,\ldots,19$.
Their masses are at least $22,22,17$. Thus
\[
 |G|\geq17+22+22+17+17=95.
\]
If the only internal singleton is at $4$, the left cap has size at
least $20$. In the right cap use its boundary radius-three ball,
the whole intervals $8,\ldots,13$ and $14,\ldots,19$, and the terminal
radius-two ball. Their supports are disjoint, so the right cap has size
at least $15+22+22+17=76$. The caps share just their boundary, giving
$|G|\geq20+76-1=95$. For the sole singleton $18$, the left cap instead
contains supports $0,\ldots,2$, $3,\ldots,8$, $9,\ldots,14$ and
$15,\ldots,18$, with sizes at least $17,22,22,15$. The same count
applies without reversing a predecessor axiom.

It remains to treat $\{4,18\}$. Let $H$ be the section on original
indices $4,\ldots,18$, with boundaries $w,z$. Their distance in $H$ is
exactly fourteen: an original root geodesic gives such a path, and
indices cannot change faster than one per edge. Thus the radius-four
hypotheses hold at both ends. Put $L_r=|\B_r^H(w)|$ and
$R_r=|\B_r^H(z)|$. Two separate disjoint packings give
\begin{align*}
 |H|&\geq L_4+22+R_3,
 &&\text{supports }[4,8],[9,14],[15,18],\\
 |H|&\geq L_3+22+R_4,
 &&\text{supports }[4,7],[8,13],[14,18].
\end{align*}
Each middle interval is a complete original nonsingleton six-layer
interval. The two packings may overlap each other; only the sets
within each one must be disjoint. Adding the inequalities and using
Lemma~\ref{lem:mixed} gives
\[
 2|H|\geq(L_3+L_4)+(R_3+R_4)+44\geq114,
 \qquad |H|\geq57.
\]
The caps outside $H$ each have size at least twenty and meet it only
at their respective boundaries. Consequently
\[
 |G|\geq20+57+20-2=95.
\]
This exhausts the alternatives and proves Theorem~\ref{thm:main}.

\section{Equality and reproducibility}

The included graph6 file
\note{data/girth_diameter_bounds/n95_k4_g5_d22.g6} contains one graph.
The property predicate reused from the prior independent positive check
computes connectedness and the exact tuple
\[
 (|V|,\delta,\Delta,\text{girth},\text{diameter},
   \text{bipartite indicator})=(95,4,4,5,22,0).
\]
It uses NetworkX~3.6.1. The portable entrypoint compares the predicate's
syntax tree to the included accepted reference before applying it.
This positive witness and Theorem~\ref{thm:main} prove the following.
\begin{corollary}
$n(4;5,22)=95$, and there is at least one such minimum graph.
\end{corollary}

The portable replay authenticates all files, verifies the $84$-pattern
cover, reconstructs all $27$ formulas and primary-variable maps, rebuilds
the unchanged RUP checker, checks every saved refutation, and computes
the positive graph's properties. The handwritten premises remain
mathematical proofs to be read; the replay is not a proof-assistant
formalization. No old graph-search cache, graph generator, SAT solver,
or private filesystem path is needed. Instructions and exact software
premises are in the companion README. The result concerns this source
row only; it neither completes the other source rows nor enumerates the
minimum graphs.

\begin{thebibliography}{9}
\bibitem{source} S. Cambie, J. Goedgebeur, J. Jooken and T. Van den Eede,
\emph{On the order-diameter ratio of girth-diameter cages},
arXiv:2511.21144v3, 19 August 2026; DMTCS 28:4, SOFSEM2026 \#1,
Table~5, physical PDF page~22.
\end{thebibliography}
\end{document}
